\insertImage{775323db-ad9f-460e-baa7-57c944b1bc0d.png}{Aquarelle d'un groupe de personnes accueillant à bras ouverts un nouvel arrivant dans leur équipe.}

\chapter{La gestion des ressources humaines}

Dans une entreprise classique, la fonction RH a un déroulé connu : elle recrute sur fiche de poste, forme sur catalogue, évalue sur grille et arbitre les conflits en dernier recours. On peut critiquer ce système — je ne m'en prive pas — mais reconnaissons-lui une qualité : chacun sait à quoi s'attendre.

Maintenant, retirez la hiérarchie. Qui recrute, quand il n'y a plus de « chef d'équipe » pour valider ? Qui décide qu'un salarié doit monter en compétences, quand personne n'est chargé de son évaluation ? Et surtout — on y viendra, c'est le sujet qui fâche — qui intervient quand deux personnes ne se supportent plus ?

Une entreprise libérée ne supprime pas les questions RH. Elle les redistribue. Et cette redistribution est l'un des chantiers les plus délicats de toute la transformation, parce qu'elle touche à ce qu'il y a de plus sensible : les personnes.

\section{Le recrutement et l'intégration des nouveaux employés}

Commençons par une évidence qu'on oublie systématiquement : dans une entreprise libérée, une erreur de recrutement coûte plus cher.

Dans une organisation classique, un recrutement moyen est amorti par la structure — le manager encadre, le process rattrape. Dans une entreprise libérée, la personne arrive dans un environnement qui suppose l'autonomie, l'initiative et l'adhésion à des valeurs fortes. Si le fit n'y est pas, il n'y a pas de manager pour compenser : c'est toute l'équipe qui absorbe.

C'est pourquoi les entreprises libérées confient très souvent le recrutement… aux équipes elles-mêmes. Chez Valve, les candidats sont évalués par leurs futurs pairs, autant sur leur capacité à travailler sans supervision que sur leurs compétences techniques\footnote{Valve Corporation (2012). Handbook for New Employees.}. La logique est imparable : qui est mieux placé pour choisir un coéquipier que l'équipe qui vivra avec ?

Mais recruter la bonne personne ne suffit pas — il faut réussir son atterrissage. Arriver dans une entreprise sans organigramme est une expérience déroutante : pas de chef à qui se référer, pas de parcours balisé, des codes implicites partout. Les organisations qui réussissent l'intégration l'ont comprise comme une transmission de culture, pas comme une formalité administrative. Chez Buurtzorg, l'organisation néerlandaise de soins à domicile, chaque nouvelle infirmière est formée à la philosophie de l'entreprise et au fonctionnement en équipe autogérée avant toute chose\footnote{Laloux, F. (2014). Reinventing Organizations: A Guide to Creating Organizations Inspired by the Next Stage of Human Consciousness. Nelson Parker.}. On ne lui apprend pas seulement son métier — elle le connaît déjà. On lui apprend comment on le pratique \emph{ici}.

Et l'intégration ne s'arrête pas au premier mois. Sans manager pour surveiller la trajectoire, c'est à l'équipe de vérifier régulièrement que le nouveau trouve sa place — et au nouveau d'oser demander de l'aide. Ce qui, dans une culture qui valorise l'autonomie, est moins facile qu'il n'y paraît.

% [KEVIN : comment recrutez-vous chez toi ? Un exemple de recrutement raté ou réussi à cause du "fit" vaudrait tous les paragraphes du monde.]

\section{La formation et le développement des compétences}

Dans une entreprise libérée, personne ne viendra vous inscrire à une formation. C'est à la fois la meilleure et la pire nouvelle du modèle.

La meilleure, parce que le développement autodirigé est puissant : chacun construit son parcours selon ses besoins réels — un cours en ligne, une conférence, un mentor en interne, un projet hors de sa zone de confort. Chez Morning Star, le transformateur de tomates californien souvent cité en exemple, chaque salarié est responsable de son propre développement et l'entreprise fournit les ressources\footnote{Laloux, F. (2014). Reinventing Organizations: A Guide to Creating Organizations Inspired by the Next Stage of Human Consciousness. Nelson Parker.}. Pour les profils curieux et autonomes, c'est un accélérateur formidable.

La pire, parce que ce système a un angle mort : ceux qui ne demandent rien. Dans le modèle classique, l'entretien annuel — malgré tous ses défauts — force au moins une conversation par an sur les compétences. Supprimez-le sans le remplacer, et les plus discrets peuvent passer des années sans que personne ne se penche sur leur progression. L'autodirection profite d'abord à ceux qui savent déjà se diriger. Les autres décrochent en silence.

La responsabilité de l'entreprise ne disparaît donc pas — elle se déplace. Il ne s'agit plus de prescrire des formations, mais de créer un environnement où se former est facile, visible et valorisé : du temps sanctuarisé, des ressources accessibles, des pairs qui partagent, et des rituels d'équipe où l'on parle de progression, pas seulement de production.

\section{La gestion des conflits et des problèmes personnels}

Voici le test décisif d'une entreprise libérée. Pas la stratégie, pas les outils : le premier vrai conflit.

Dans une organisation classique, le conflit a un chemin : le manager, puis les RH. Ce chemin est souvent imparfait, mais il existe. Dans une entreprise libérée, le premier réflexe est de croire que la bonne ambiance et les valeurs partagées suffiront. Elles ne suffisent jamais. Des adultes qui travaillent ensemble finissent par s'opposer — sur une décision, une charge de travail, une manière d'être. C'est la vie normale d'un collectif.

Sans mécanisme prévu, un conflit dans une entreprise libérée a un comportement particulier : il dure. Personne n'a le rôle de l'arbitre, tout le monde attend que « ça se règle entre eux », et la transparence ambiante fait que tout le monde assiste au spectacle. Ce qui aurait été une friction de bureau devient un feuilleton d'entreprise.

Les organisations matures s'outillent donc explicitement : des processus de médiation par les pairs, une escalade progressive et connue de tous — d'abord le face-à-face, puis un médiateur choisi ensemble, puis un panel — comme le pratique précisément Morning Star\footnote{Laloux, F. (2014). Reinventing Organizations: A Guide to Creating Organizations Inspired by the Next Stage of Human Consciousness. Nelson Parker.}. Chez Buurtzorg, des coachs d'équipe — qui ne sont pas des managers — aident les équipes à traverser leurs tensions sans décider à leur place. Le point commun : le chemin du conflit est défini \emph{avant} le conflit. Après, il est trop tard.

Reste les problèmes personnels — la santé, la famille, les coups durs. Une entreprise qui demande à ses salariés de s'engager en tant que personnes entières doit accepter de les soutenir en tant que personnes entières. C'est la contrepartie du modèle, et elle n'est pas optionnelle : on ne peut pas vouloir l'engagement du lundi et ignorer la détresse du jeudi.

Une entreprise libérée ne se juge pas à ses journées faciles. Elle se juge à la façon dont elle traite ses moments difficiles.
